% =============================================
% CHƯƠNG 4 — PHÂN TÍCH VÀ THIẾT KẾ HỆ THỐNG (13 trang) ★ LÕI ĐỒ ÁN
% Status: OUTLINE — chờ draft từ agent
% Đặc điểm: 90% tài sản đã có sẵn, chủ yếu biên tập
% =============================================

\chapter{Phân tích và thiết kế hệ thống}
\label{ch:thiet-ke}

% [INTRO CHƯƠNG — 1 đoạn]
% Chương này trình bày thiết kế tổng thể và chi tiết của hệ thống, bao gồm:
% kiến trúc 3 tầng, lược đồ dữ liệu, luồng xác thực, thuật toán QR HMAC xoay,
% cơ chế xác minh ngang hàng, hệ thống chống gian lận và tính năng cổng máy chiếu.

% ─────────────────────────────────────────────
\section{Kiến trúc tổng thể}
\label{sec:kien-truc}
% Ước lượng: 1.5 trang
% Nguồn: Notion page 1 + docs/architecture-diagram.html
% Nội dung:
%   §4.1.1 Mô hình 3 tầng: Client (Mini App + Admin Web) → Backend (Firebase) → Storage
%   §4.1.2 Phân tách trách nhiệm:
%     - Client: UI, validation, capture face/QR
%     - Backend: validate, HMAC sign/verify, fraud detection
%     - Storage: persist + realtime sync
%   §4.1.3 Component diagram tổng quát
%   Hình 4.1: Sơ đồ kiến trúc 3-tier (export từ Notion mermaid)
%   Hình 4.2: Component diagram

% ─────────────────────────────────────────────
\section{Thiết kế cơ sở dữ liệu Firestore}
\label{sec:db-design}
% Ước lượng: 2 trang
% Nguồn: src/types/index.ts (16 interfaces) + firestore.rules + Notion page 6
% Nội dung:
%   §4.2.1 Mô hình dữ liệu: document/collection vs relational
%   §4.2.2 12 collection chính:
%     users, classes, sessions, sessions/secrets, attendance,
%     face_registrations, verified_students, oauth_states,
%     fraud_reports, absence_requests, teacher_invites, pairing_tokens
%   §4.2.3 Schema từng collection:
%     Bảng 4.1: users (fields, types, indexes)
%     Bảng 4.2: classes
%     Bảng 4.3: sessions + sessions/secrets (subcollection)
%     Bảng 4.4: attendance
%     ... (5-6 bảng chính, các collection nhỏ gộp lại)
%   §4.2.4 ER diagram (Notion page 6 → PNG)
%   Hình 4.3: Sơ đồ thực thể-liên kết Firestore
%   §4.2.5 Lý do chọn Firestore (denormalization, security rules, realtime)

% ─────────────────────────────────────────────
\section{Thiết kế luồng xác thực và phân quyền}
\label{sec:auth-design}
% Ước lượng: 1.5 trang
% Nguồn: Notion page 3 + src/services/auth.service.ts
% Nội dung:
%   §4.3.1 Sequence: Zalo OAuth → Custom Token → Firebase Auth
%   §4.3.2 3 vai trò và phân quyền:
%     Bảng 4.5 — Ma trận quyền theo collection × role
%   §4.3.3 Xác minh email HUST (sinh viên):
%     - sendOTP → verifyOTP → setDoc(verified_students)
%     - Bypass cho tài khoản demo (MSSV 20225413)
%   §4.3.4 Mã mời giảng viên (teacher_invites)
%   §4.3.5 Firestore Security Rules (snippet 5-7 dòng minh hoạ)
%   Hình 4.4: Sequence diagram đăng nhập (Notion page 3)
%   Hình 4.5: State diagram của user

% ─────────────────────────────────────────────
\section{Thuật toán QR HMAC xoay 30 giây}
\label{sec:hmac-rotating}
% Ước lượng: 1.5 trang
% Nguồn: src/utils/crypto.ts + src/hooks/useQRGenerator.ts + báo cáo 11-16/05
% Nội dung:
%   §4.4.1 Vấn đề: QR tĩnh dễ bị screenshot và share
%   §4.4.2 Giải pháp: QR động xoay 30s với HMAC signature
%   §4.4.3 Thuật toán tạo QR (pseudocode):
%     INPUT: sessionId, secret, anchor (session.startedAt), refreshInterval=30s
%     periodIndex = floor((now - anchor) / refreshInterval)
%     timestamp = anchor + periodIndex * refreshInterval
%     nonce = randomBytes(8)
%     payload = {sessionId, timestamp, nonce}
%     signature = HMAC-SHA256(secret, JSON(payload))
%     qrContent = base64(payload + signature)
%   §4.4.4 Thuật toán xác minh phía SV
%   §4.4.5 Đồng bộ wall-clock giữa GV (Mini App) và máy chiếu (Web Admin)
%     - Dùng cùng anchor → cùng periodIndex → cùng thời điểm refresh
%   §4.4.6 Phân tích an toàn: chống screenshot share (max 30s + 90s tolerance)
%   Hình 4.6: Sơ đồ trạng thái QR xoay theo thời gian
%   Hình 4.7: Pseudocode minh hoạ

% ─────────────────────────────────────────────
\section{Thiết kế UI/UX và bộ nhận diện}
\label{sec:ui-design}
% Ước lượng: 1.5 trang
% Nguồn: CLAUDE.md Design System + báo cáo 11-16/05 (UI Bách Khoa)
% Nội dung:
%   §4.5.1 Bộ nhận diện Bách Khoa: đỏ HUST #be1d2c, logo SVG, ảnh tòa C1
%   §4.5.2 Design tokens: màu sắc, typography, spacing, shadow
%   §4.5.3 Component library tự xây: ScoreRing, DarkStatCard, DarkModal, ...
%   §4.5.4 Pattern: header HUST đỏ, card trắng + shadow, bottom-sheet modal
%   §4.5.5 Light theme chính, accent đỏ HUST
%   Hình 4.8: Mock-up một số màn hình chính (đã capture)

% ─────────────────────────────────────────────
\section{Luồng điểm danh 4 bước cho sinh viên}
\label{sec:attendance-flow}
% Ước lượng: 2 trang
% Nguồn: Notion page 4 + src/pages/student/StudentAttendance.tsx
% Nội dung:
%   §4.6.1 Tổng quan 4 bước
%   §4.6.2 Bước 1 — Quét QR động của GV:
%     - InlineQRScanner + jsQR
%     - Validate HMAC + timestamp (90s expiry)
%   §4.6.3 Bước 2 — Xác minh khuôn mặt:
%     - Auto-scan, oval clip, ngưỡng 70%
%   §4.6.4 Bước 3 — Trao đổi QR với bạn xung quanh (peer-to-peer):
%     - Side-by-side: hiện QR riêng + scan QR bạn
%     - Tối thiểu 3 peers
%   §4.6.5 Bước 4 — Hoàn tất + trust score
%   Hình 4.9: Flowchart 4 bước (Notion page 4)
%   Hình 4.10: Screenshot từng bước

% ─────────────────────────────────────────────
\section{Hệ thống chống gian lận 5 lớp}
\label{sec:anti-fraud}
% Ước lượng: 2 trang
% Nguồn: Notion page 4 + functions/src/services/fraud.service.ts
% Nội dung:
%   §4.7.1 Tổng quan: defense in depth
%   §4.7.2 Lớp 1 — QR HMAC xoay (đã trình bày §4.4)
%   §4.7.3 Lớp 2 — Face verification (ngưỡng 70%)
%   §4.7.4 Lớp 3 — Peer verification (≥3 peers)
%   §4.7.5 Lớp 4 — GPS geofencing (phạm vi 200m)
%   §4.7.6 Lớp 5 — Trust score & Fraud Detection:
%     - Cloud Function quét pattern (high-score buddy, GPS-mismatch, time-anomaly, ...)
%     - 7 loại fraud type
%   §4.7.7 Workflow xử lý fraud report (GV review → SV kháng nghị)
%   Hình 4.11: Sơ đồ kiến trúc anti-fraud 5 lớp
%   Bảng 4.6: Phân loại 7 loại gian lận và cách phát hiện

% ─────────────────────────────────────────────
\section{Tính năng Cổng máy chiếu (Teacher Display)}
\label{sec:projector}
% Ước lượng: 1.5 trang
% Nguồn: memory feature-projector-pairing.md + báo cáo 11-16/05
% Nội dung:
%   §4.8.1 Bối cảnh: QR điện thoại quá nhỏ cho lớp đông
%   §4.8.2 Kiến trúc pairing:
%     - Web tạo token 128-bit hết hạn 60s
%     - GV scan token bằng Mini App → claim transaction
%   §4.8.3 Firestore one-shot rule: pending → paired (không revert, không claim 2 lần)
%   §4.8.4 Đồng bộ countdown wall-clock giữa Mini App ↔ Web
%   §4.8.5 Wake Lock API chống máy chiếu sleep
%   §4.8.6 3 trạng thái UI: Pairing → Displaying → Ended
%   Hình 4.12: Sequence diagram pairing flow
%   Hình 4.13: 3 màn hình UI cổng chiếu

% ─────────────────────────────────────────────
% \section*{Kết chương}
% Đã trình bày kiến trúc và 8 thành phần thiết kế cốt lõi. Chương 5 trình bày
% triển khai cụ thể và kết quả kiểm thử.
